% !TeX root = ../../Book3.tex
% 纲要标题：### 22.5 整数流闭包与紧性的证明
% 纲要位置：几何学计划纲要.md，第 2892 行。
% 正文状态：未撰写；本文件目前是框架。

% 纲要内容开始
% 
% * 正规流的弱紧性
% 
%   在共同紧支和质量、边界质量界下，先获得对测试形式收敛的子列。
% 
% * 变形估计与统一平坦控制
% 
%   利用多面体逼近把弱序列转化为可控制的几何误差。
% 
% * 切片及维数归纳
% 
%   证明几乎所有切片保留相容的整数信息，逐维恢复整性。
% 
% * 极限的可整流性
% 
%   此处单列闭包定理真正需要的切向结构论证，不从弱收敛直接断言。
% 
% * 整数重数的保持
% 
%   区分可整流性与重数属于整数这两项结论。
% 
% * 边界相容及下半连续
% 
%   证明极限边界正确，并核验质量下半连续。
% 
% * 紧性定理及非紧推广的边界
% 
%   先完成固定紧支版本，再说明局部质量界与质量逃逸的区别。
% 
% ---
% 
% 纲要内容结束

\section{整数流闭包与紧性的证明}
\label{g3:sec:2205}

待撰写

\subsection{正规流的弱紧性}
\label{g3:subsec:2205:01}
待撰写

\subsection{变形估计与统一平坦控制}
\label{g3:subsec:2205:02}
待撰写

\subsection{切片及维数归纳}
\label{g3:subsec:2205:03}
待撰写

\subsection{极限的可整流性}
\label{g3:subsec:2205:04}
待撰写

\subsection{整数重数的保持}
\label{g3:subsec:2205:05}
待撰写

\subsection{边界相容及下半连续}
\label{g3:subsec:2205:06}
待撰写

\subsection{紧性定理及非紧推广的边界}
\label{g3:subsec:2205:07}
待撰写
% 纲要条目排为编号小节标题，并进入目录与书签；说明段仍保留于纲要注释中。
